\section{Задание}

Моделируем информационный центр. В информационный центр приходят клиенты (пользователи) через интервал времени 10 ± 2 минуты. Если все три имеющихся оператора заняты, клиенту отказывают в обслуживании. Операторы имеют разную производительность и могут обеспечивать обслуживание среднего запроса от пользователя за 20 ± 5, 40 ± 10 и 40 ± 20 ед. времени (минут). Клиенты стараются занять свободного оператора с максимальной производительностью. Полученные запросы сдаются в накопитель, откуда выбираются на обработку. На первый компьютер — от первого и второго операторов, на второй — от третьего. Время обработки запроса в компьютерах — 15 и 30 минут соответственно. Смоделировать процесс обработки 300 запросов. Определить вероятность отказа.

\section{Теоретическая часть}

\subsection{Схемы модели}

На рисунке 1 представлена структурная схема модели.

\includeimage
	{blockDiagram}{f}{H}{0.8\textwidth}
	{Структурная схема модели}

В процессе взаимодействия клиентов с информационным центром возможно два режима работы:

\begin{itemize}
    \item режим нормального обслуживания, когда клиент выбирает одного из свободных операторов, отдавая предпочтение тому, у кого максимальная производительность;
    \item режим отказа клиенту в обслуживании, когда все операторы заняты.
\end{itemize}

На рисунке 2 представлена схема модели в терминах систем массового обслуживания (СМО).

\includeimage
	{queuingSystems}{f}{H}{0.7\textwidth}
	{Схема модели в терминах СМО}

\subsection{Событийный принцип моделирования}

Событийный принцип использует дискретный характер изменения состояний моделируемой системы. Состояния блоков имитационной модели анализируются только в момент появления события. Момент наступления следующего события определяется минимальным значением из списка будущих событий, представляющего совокупность ближайших изменений состояния каждого блока системы.

Этот подход более эффективен по сравнению с пошаговым, так как устраняет необходимость проверки состояния системы на каждом временном шаге.

В разработанной модели определены следующие типы событий:

\begin{itemize}
    \item \textbf{CLIENT\_ARRIVAL} — прибытие клиента в информационный центр
    \item \textbf{OPERATOR\_FINISH} — окончание обработки запроса оператором
    \item \textbf{COMPUTER\_FINISH} — окончание обработки запроса компьютером
\end{itemize}

\subsection{Равномерное распределение}

Случайная величина $X$ имеет равномерное распределение на отрезке $[a, b]$, если её плотность распределения $f(x)$ равна:

\begin{equation}
p(x) = \begin{cases}
\frac{1}{b-a}, & \text{если } a \leq x \leq b \\
0, & \text{иначе}
\end{cases}
\end{equation}

При этом функция распределения $F(x)$ равна:

\begin{equation}
F(x) = \begin{cases}
0, & x < a \\
\frac{x-a}{b-a}, & a \leq x \leq b \\
1, & x > b
\end{cases}
\end{equation}

Обозначение: $X \sim R[a, b]$.

Интервал времени между появлением $i$-й и $(i-1)$-й заявки по равномерному закону вычисляется следующим образом:

\begin{equation}
T_i = a + (b - a) \cdot R
\end{equation}

где $R$ — псевдослучайное число из интервала $[0, 1]$.

\subsection{Переменные и уравнение имитационной модели}

\textbf{Эндогенные переменные:}

\begin{itemize}
    \item время обработки задания $i$-ым оператором;
    \item время решения задания на $j$-ом компьютере.
\end{itemize}

\textbf{Экзогенные переменные:}

\begin{itemize}
    \item $n_0$ — число обслуженных клиентов;
    \item $n_1$ — число клиентов, получивших отказ.
\end{itemize}

Вероятность отказа в обслуживании клиента вычисляется как:

\begin{equation}
P = \frac{n_1}{n_0 + n_1}
\end{equation}

\section{Описание реализации}

\subsection{Используемые библиотеки}

Программа реализована на языке Python с использованием следующих библиотек:

\begin{itemize}
    \item \textbf{PyQt6} — создание графического интерфейса пользователя
    \item \textbf{random} — генерация псевдослучайных чисел
    \item \textbf{typing} — аннотации типов для улучшения читаемости кода
    \item \textbf{enum} — определение перечислений для типов событий
\end{itemize}

\subsection{Структура программы}

Программа состоит из следующих модулей:

\begin{itemize}
    \item \texttt{main.py} — точка входа в приложение
    \item \texttt{gui/main\_window.py} — главное окно с параметрами моделирования
    \item \texttt{gui/connections\_window.py} — окно настройки связей между элементами
    \item \texttt{gui/results\_window.py} — окно отображения результатов
    \item \texttt{models/distributions.py} — классы для генерации случайных величин
    \item \texttt{models/elements.py} — классы элементов системы (генератор, операторы, компьютеры, очереди)
    \item \texttt{models/event\_model.py} — реализация событийного подхода моделирования
    \item \texttt{constants.py} — константы интерфейса и параметры по умолчанию
\end{itemize}

\subsection{Метод решения}

Моделирование информационного центра выполняется с использованием событийного подхода:

\begin{enumerate}
    \item Инициализация списка будущих событий с первым событием прибытия клиента
    \item Основной цикл моделирования до обработки заданного числа клиентов:
    \begin{itemize}
        \item Выбор ближайшего события из отсортированного списка
        \item Обновление текущего времени моделирования
        \item Обработка события в зависимости от его типа:
        \begin{itemize}
            \item \textbf{Прибытие клиента}: поиск свободного оператора с максимальной производительностью или отказ в обслуживании
            \item \textbf{Окончание обработки оператором}: добавление запроса в очередь компьютера и запуск обработки при возможности
            \item \textbf{Окончание обработки компьютером}: извлечение следующего запроса из очереди при наличии
        \end{itemize}
        \item Добавление новых событий в список с сохранением сортировки по времени
    \end{itemize}
    \item Сбор статистики по очередям (максимальный и средний размер)
\end{enumerate}

Производительность оператора вычисляется как величина, обратная среднему времени обработки:

\begin{equation}
\text{productivity} = \frac{1}{\text{mean\_processing\_time}}
\end{equation}

Клиент выбирает свободного оператора с наибольшей производительностью.

\section{Практическая часть}

На рисунке 3 представлен результат работы разработанной программы для нахождения вероятности отказа с параметрами по умолчанию:

\begin{itemize}
    \item Генератор клиентов: равномерное распределение с параметрами $a = 8$, $b = 12$ (среднее 10 ± 2 минуты)
    \item Оператор 1: равномерное распределение с параметрами $a = 15$, $b = 25$ (среднее 20 ± 5 минут)
    \item Оператор 2: равномерное распределение с параметрами $a = 30$, $b = 50$ (среднее 40 ± 10 минут)
    \item Оператор 3: равномерное распределение с параметрами $a = 20$, $b = 60$ (среднее 40 ± 20 минут)
    \item Компьютер 1: детерминированное время обработки 15 минут
    \item Компьютер 2: детерминированное время обработки 30 минут
    \item Количество клиентов: 300
\end{itemize}

\includeimage
	{result}{f}{H}{1.0\textwidth}
	{Результат работы программы}

Программа корректно вычисляет вероятность отказа в обслуживании клиентов. Событийный подход обеспечивает точное моделирование без погрешностей, связанных с дискретизацией времени.

\section{Расширенный функционал}

Помимо базовой функциональности, реализованной в соответствии с заданием, программа поддерживает дополнительные возможности:

\begin{itemize}
    \item \textbf{Динамическое количество операторов и компьютеров} — пользователь может добавлять и удалять операторы и компьютеры, настраивая конфигурацию системы под различные сценарии моделирования
    \item \textbf{Настройка связей между элементами} — специальное окно позволяет визуально настраивать, какие операторы отправляют запросы на какие компьютеры
    \item \textbf{Статистика по очередям} — программа отслеживает максимальный и средний размер каждой очереди в процессе моделирования
\end{itemize}

Данный расширенный функционал позволяет исследовать различные конфигурации информационного центра и проводить сравнительный анализ эффективности разных топологий системы.

\section{Вывод}

В ходе выполнения лабораторной работы была разработана программа с графическим интерфейсом для моделирования работы информационного центра с использованием событийного принципа.

Программа обеспечивает гибкую настройку параметров генератора клиентов, операторов и компьютеров, автоматическую валидацию входных данных и наглядное представление результатов моделирования. Реализованный событийный подход обеспечивает точное моделирование без погрешностей дискретизации, характерных для пошагового метода.

Расширенный функционал программы позволяет исследовать различные конфигурации информационного центра с произвольным количеством операторов и компьютеров, что делает её удобным инструментом для анализа систем массового обслуживания.
